\chapter*{Notazione e convenzioni}
\addcontentsline{toc}{chapter}{Notazione e convenzioni}
\markboth{Notazione e convenzioni}{Notazione e convenzioni}

I simboli seguenti hanno lo stesso significato in tutta la tesi. Dove un simbolo
corrisponde a un campo di configurazione del codice, il nome del campo è riportato
in carattere \code{monospaziato}: è quel nome, e non il simbolo, a comparire nei
file \file{config.json} salvati da ogni esecuzione.

\section*{Forme tensoriali}

Ogni tensore è indicato con la sua forma fra parentesi tonde, dalla dimensione più
esterna alla più interna: $\shape{\dbatch, \dseq, \demb}$ denota un tensore di
$\dbatch$ elementi di batch, ciascuno una sequenza di $\dseq$ posizioni, ciascuna
un vettore di $\demb$ componenti. Le trasformazioni sono scritte come catene di
forme; una catena incompleta è un errore, non un'abbreviazione.

\section*{Transformer autoregressivo (capitolo 3)}

\begin{center}
\begin{tabularx}{\textwidth}{@{}l l X@{}}
\toprule
\textbf{Simbolo} & \textbf{Campo} & \textbf{Significato} \\
\midrule
$\dbatch$   & \code{batch\_size} & Numero di sequenze per batch \\
$\dseq$     & \code{block\_size} & Lunghezza di contesto \\
$\demb$     & \code{n\_embd}     & Dimensione di embedding del modello \\
$\nhead$    & \code{n\_head}     & Numero di teste di attenzione \\
$\dhead$    & \code{head\_size}  & Dimensione per testa, $\dhead = \demb/\nhead$ \\
$\dvoc$     & \code{vocab\_size} & Cardinalità del vocabolario del tokenizer \\
$\mathcal{L}$ & ---              & Cross-entropy media per token \\
$\PPL$      & ---                & Perplexity, $\PPL = \exp(\mathcal{L})$ \\
\bottomrule
\end{tabularx}
\end{center}

\section*{Vision Transformer (capitoli 4 e 7)}

\begin{center}
\begin{tabularx}{\textwidth}{@{}l l X@{}}
\toprule
\textbf{Simbolo} & \textbf{Campo} & \textbf{Significato} \\
\midrule
$\nch$      & \code{n\_channels} & Canali dell'immagine di ingresso \\
$H, W$      & \code{image\_size} & Altezza e larghezza dell'immagine \\
$p$         & \code{patch\_size} & Lato della patch quadrata \\
$\npatch$   & \code{n\_patches}  & Numero di patch, $\npatch = (H/p)^2$ \\
$\dpatch$   & \code{patch\_dim}  & Patch appiattita, $\dpatch = p^2\nch$ \\
$\dvit$     & \code{embed\_dim}  & Dimensione dei token del ViT \\
--- & \code{seq\_len} & Lunghezza di sequenza, $\npatch + 1$ (token \code{[CLS]}) \\
\bottomrule
\end{tabularx}
\end{center}

\section*{Circuito quantistico variazionale}

\begin{center}
\begin{tabularx}{\textwidth}{@{}l l X@{}}
\toprule
\textbf{Simbolo} & \textbf{Campo} & \textbf{Significato} \\
\midrule
$\nq$ & \code{n\_qubits}  & Numero di qubit simulati \\
$\nqlayers$ & \code{n\_qlayers} & Profondità di \code{StronglyEntanglingLayers} \\
$\ket{\psi}$ & --- & Stato di $\nq$ qubit, vettore in $\C^{2^{\nq}}$ \\
$\Uenc(x)$ & --- & Unitaria di codifica dei dati \\
$\expval{Z_i}$ & --- & Valore di aspettazione di Pauli-$Z$ sul qubit $i$, in $[-1,1]$ \\
$W$ & \code{weights} & Pesi variazionali, forma $\shape{\nqlayers, \nq, 3}$ \\
\bottomrule
\end{tabularx}
\end{center}

\section*{Metriche e statistica (capitoli 6 e 8)}

\begin{center}
\begin{tabularx}{\textwidth}{@{}l X@{}}
\toprule
\textbf{Simbolo} & \textbf{Significato} \\
\midrule
$\acctr, \accte$ & Accuracy su training set e su test set \\
$\gap$ & Gap di generalizzazione, $\gap = \acctr - \accte$ \\
$d_s$ & Differenza appaiata sul seed $s$: $d_s = m_{q,s} - m_{c,s}$ \\
$\dmean, \sd$ & Media e deviazione standard campionaria delle differenze \\
$\dz$ & Dimensione dell'effetto appaiata, $\dz = \dmean / \sd$ \\
$\IC$ & Intervallo di confidenza al \SI{95}{\percent} \\
$\alpha$ & Livello di significatività, fissato a $0{,}05$ \\
\bottomrule
\end{tabularx}
\end{center}

\section*{Convenzione di segno}

\convsegno{} Il pedice $q$ indica la variante quantistica, il pedice $c$ la
baseline classica appaiata. La convenzione è ripetuta nella didascalia di ogni
tabella del capitolo~8 perché è la principale fonte di errore in lettura: una
differenza negativa sull'\emph{accuracy} significa che il modello quantistico è
peggiore, mentre una differenza negativa sul \emph{gap} significa che generalizza
meglio.

\section*{Registri del testo}

Per la regola redazionale adottata, i quattro registri del discorso sono
tipograficamente distinti dove il rischio di confusione è concreto:

\begin{definizione}[Esempio]
Enunciato teorico, valido indipendentemente da questa implementazione.
\end{definizione}

\begin{implementazione}
Ciò che il codice fa realmente, con riferimento a classe, funzione e file.
\end{implementazione}

\begin{osservazione}
Un valore misurato, con l'indicazione della run da cui proviene.
\end{osservazione}

\begin{interpretazione}
Una spiegazione proposta, con l'indicazione esplicita di quanto sia sostenuta
dai dati e quanto resti congetturale.
\end{interpretazione}

\cleardoublepage
